\documentclass[9pt,a4paper]{extarticle}

% ==========================================
% Basic packages
% ==========================================
\usepackage[T1]{fontenc}
\usepackage[utf8]{inputenc}
\usepackage[brazil]{babel}
\usepackage{geometry}
\geometry{margin=2cm}
\usepackage{parskip}
\usepackage{microtype}

% Math
\usepackage{amsmath, amssymb, amsthm, bm}

% ==========================================
% Document
% ==========================================
\begin{document}

\section*{Lecture Notes: Newton's Method}

\section{Motivation}

Newton's method is a second-order optimization technique that uses information
from the gradient and the Hessian (matrix of second derivatives) of the objective
function to construct a local quadratic model around the current iterate. This
quadratic approximation is then minimized to obtain a new point.

The basic idea is: if the function $f:\mathbb{R}^n\to\mathbb{R}$ is smooth and
we are at a point $x_i$ close to a minimizer, then a quadratic Taylor
approximation of $f$ around $x_i$ can be used to compute a search direction that
points directly toward the optimum (for quadratic functions, in a single step).

\section{Quadratic Approximation and Newton Step}

\subsection*{Step 1: Taylor expansion around the current point}

Consider a twice continuously differentiable function $f:\mathbb{R}^n \to \mathbb{R}$.
Around a point $x_i$, we can approximate $f$ by a second-order Taylor expansion:
\[
f(x) \approx f(x_i)
+ \nabla f_i^\top (x - x_i)
+ \tfrac12 (x - x_i)^\top H_i (x - x_i),
\]
where
\[
\nabla f_i := \nabla f(x_i)
\quad\text{e}\quad
H_i := \nabla^2 f(x_i)
\]
é a matriz Hessiana avaliada em $x_i$.

\subsection*{Step 2: Local quadratic model}

Definimos o modelo quadrático local $m_i:\mathbb{R}^n\to\mathbb{R}$ como
\[
m_i(x)
= f(x_i)
+ \nabla f_i^\top (x - x_i)
+ \tfrac12 (x - x_i)^\top H_i (x - x_i).
\]

Queremos encontrar o ponto que minimiza esse modelo. Seja
\[
p := x - x_i
\]
o deslocamento a partir de $x_i$. Então
\[
m_i(x_i + p)
= f(x_i) + \nabla f_i^\top p + \tfrac12 p^\top H_i p.
\]

O problema local de minimização torna-se:
\[
\min_{p \in \mathbb{R}^n} \;
q(p) := \nabla f_i^\top p + \tfrac12 p^\top H_i p.
\]

\subsection*{Step 3: Condição de otimalidade do modelo quadrático}

Para minimizar $q(p)$, derivamos em relação a $p$ e igualamos a zero:
\[
\nabla_p q(p)
= \nabla f_i + H_i p = 0.
\]

Assim, a condição de primeira ordem para o mínimo do modelo é
\[
H_i p = - \nabla f_i.
\]

Supondo que $H_i$ seja não singular (invertível), obtemos
\[
p^* = - H_i^{-1} \nabla f_i.
\]

\subsection*{Step 4: Atualização de Newton}

O próximo ponto $x_{i+1}$ é obtido somando o deslocamento ótimo $p^*$ ao ponto atual:
\[
x_{i+1} = x_i + p^*
= x_i - H_i^{-1} \nabla f_i.
\]

Esta é a fórmula clássica do passo de Newton:
\[
\boxed{
x_{i+1} = x_i - \bigl[\nabla^2 f(x_i)\bigr]^{-1} \nabla f(x_i).
}
\]

Como a aproximação de Taylor truncou termos de ordem superior, essa atualização
é aplicada de forma iterativa para aproximar um ponto $x^*$ tal que $\nabla f(x^*)=0$.

\section{Geometric Interpretation}

Em uma dimensão ($n=1$), Newton para encontrar raízes de $f'(x)=0$ interpreta-se
como:
\begin{itemize}
    \item A reta tangente ao gráfico de $f'(x)$ em $x_i$ é usada para encontrar
          o próximo ponto onde a derivada se anula.
    \item Em múltiplas dimensões, o mesmo raciocínio é generalizado usando o
          gradiente $\nabla f$ e a Hessiana $\nabla^2 f$ para determinar o
          deslocamento $p^*$ que anula o gradiente do modelo quadrático.
\end{itemize}

Quando $x_i$ está suficientemente próximo de um mínimo estrito e $H_i$ é
definida positiva, o método apresenta convergência quadrática.

\section{Newton's Method Algorithm (Basic Version)}

We summarize the basic Newton method for unconstrained minimization.

\begin{enumerate}
    \item Choose an initial guess $x_1$.
    \item For $i = 1,2,\dots$:
    \begin{enumerate}
        \item Compute the gradient and Hessian at $x_i$:
        \[
        \nabla f_i = \nabla f(x_i),
        \qquad
        H_i = \nabla^2 f(x_i).
        \]
        \item Solve the linear system
        \[
        H_i p_i = - \nabla f_i
        \]
        for the Newton step $p_i$.
        \item Update
        \[
        x_{i+1} = x_i + p_i.
        \]
        \item If $\|\nabla f_{i+1}\|$ is sufficiently small, stop.
    \end{enumerate}
\end{enumerate}

\section{Special Case: Quadratic Functions}

Considere uma função quadrática geral
\[
f(x) = \tfrac12 x^\top A x + b^\top x + c,
\]
onde $A$ é uma matriz simétrica definida positiva. Então:
\[
\nabla f(x) = A x + b,
\qquad
\nabla^2 f(x) = A.
\]

O ponto de mínimo é dado por
\[
\nabla f(x^*) = 0
\quad\Longrightarrow\quad
A x^* + b = 0
\quad\Longrightarrow\quad
x^* = - A^{-1} b.
\]

O passo de Newton, em qualquer iteração $i$, é
\[
x_{i+1}
= x_i - A^{-1} (A x_i + b)
= x_i - A^{-1} A x_i - A^{-1} b
= - A^{-1} b
= x^*.
\]

Ou seja, para funções quadráticas com Hessiana constante $A$, o método de
Newton encontra o ponto ótimo em \emph{uma única iteração} (independentemente
de $x_1$), desde que $A$ seja não singular.

\section{Example: Minimization of a Quadratic Function}

Considere a função
\[
f(x_1, x_2) = x_1 - x_2 + 2x_1^2 + 2x_1 x_2 + x_2^2,
\]
com ponto inicial
\[
x_1 =
\begin{bmatrix}
0\\
0
\end{bmatrix}.
\]

\subsection*{Step 1: Gradient and Hessian}

O gradiente é
\[
\nabla f(x_1, x_2) =
\begin{bmatrix}
\frac{\partial f}{\partial x_1}\\[4pt]
\frac{\partial f}{\partial x_2}
\end{bmatrix}
=
\begin{bmatrix}
1 + 4x_1 + 2x_2\\[4pt]
-1 + 2x_1 + 2x_2
\end{bmatrix}.
\]

O Hessiano (constante) é
\[
\nabla^2 f(x) =
\begin{bmatrix}
\frac{\partial^2 f}{\partial x_1^2} & \frac{\partial^2 f}{\partial x_1 \partial x_2}\\[4pt]
\frac{\partial^2 f}{\partial x_2 \partial x_1} & \frac{\partial^2 f}{\partial x_2^2}
\end{bmatrix}
=
\begin{bmatrix}
4 & 2\\[4pt]
2 & 2
\end{bmatrix}
= H.
\]

\subsection*{Step 2: Evaluate at the initial point}

No ponto inicial $x_1 = (0,0)^\top$,
\[
\nabla f_1 = \nabla f(0,0) =
\begin{bmatrix}
1\\
-1
\end{bmatrix}.
\]

\subsection*{Step 3: Compute the Newton step}

Precisamos de $H^{-1}$. Como
\[
H =
\begin{bmatrix}
4 & 2\\[4pt]
2 & 2
\end{bmatrix},
\]
uma inversa possível é
\[
H^{-1}
=
\begin{bmatrix}
\frac12 & -\frac12\\[4pt]
-\frac12 & 1
\end{bmatrix}.
\]

O passo de Newton é
\[
p_1 = - H^{-1} \nabla f_1
= -
\begin{bmatrix}
\frac12 & -\frac12\\[4pt]
-\frac12 & 1
\end{bmatrix}
\begin{bmatrix}
1\\
-1
\end{bmatrix}
=
-
\begin{bmatrix}
\frac12(1) + (-\frac12)(-1)\\[4pt]
-\frac12(1) + 1(-1)
\end{bmatrix}
=
-
\begin{bmatrix}
1\\[4pt]
-\tfrac32
\end{bmatrix}
=
\begin{bmatrix}
-1\\[4pt]
\tfrac32
\end{bmatrix}.
\]

\subsection*{Step 4: Update the point}

\[
x_2 = x_1 + p_1
=
\begin{bmatrix}
0\\
0
\end{bmatrix}
+
\begin{bmatrix}
-1\\[4pt]
\frac32
\end{bmatrix}
=
\begin{bmatrix}
-1\\[4pt]
\frac32
\end{bmatrix}.
\]

\subsection*{Step 5: Check optimality}

Calculamos o gradiente em $x_2$:
\[
\nabla f_2 = \nabla f(-1, \tfrac32)
=
\begin{bmatrix}
1 + 4(-1) + 2(\tfrac32)\\[4pt]
-1 + 2(-1) + 2(\tfrac32)
\end{bmatrix}
=
\begin{bmatrix}
1 - 4 + 3\\[4pt]
-1 - 2 + 3
\end{bmatrix}
=
\begin{bmatrix}
0\\[4pt]
0
\end{bmatrix}.
\]

Logo, $\nabla f_2 = 0$ e $x_2$ é o ponto ótimo.
Como esperado para uma função quadrática, o método de Newton convergiu em
uma única iteração.

\section{Newton's Method with Line Search}

Na prática, para funções não quadráticas, o passo
\[
x_{i+1} = x_i - H_i^{-1} \nabla f_i
\]
pode não ser adequado:
\begin{itemize}
    \item o método pode divergir;
    \item pode convergir para um ponto de sela ou máximo;
    \item o tamanho do passo pode ser muito grande ou muito pequeno.
\end{itemize}

Uma modificação comum é introduzir um \emph{line search} na direção de Newton:
\[
s_i = - H_i^{-1} \nabla f_i,
\]
\[
x_{i+1} = x_i + \lambda_i^* s_i,
\]
onde $\lambda_i^*$ é um passo que tenta minimizar $f(x_i + \lambda s_i)$ em relação a $\lambda$ (por exemplo, satisfazendo condições de Wolfe). Isso leva ao \emph{método de Newton com busca linear}, que em geral é mais robusto e evita convergência para pontos indesejáveis.

\section{Weaknesses of Newton's Method}

Apesar de suas boas propriedades teóricas (convergência quadrática próxima ao
ótimo, convergência em um passo para funções quadráticas), o método de Newton
tem limitações práticas importantes.

\subsection*{Cálculo e armazenamento da Hessiana}

O método requer a matriz Hessiana $H_i = \nabla^2 f(x_i)$:
\begin{itemize}
    \item Para problemas de dimensão $n$, $H_i$ é uma matriz $n \times n$, com
          $O(n^2)$ elementos.
    \item Em muitos problemas práticos (como em aprendizado de máquina com
          milhares ou milhões de parâmetros), é inviável calcular e armazenar
          explicitamente a Hessiana.
\end{itemize}

\subsection*{Custo de inversão / solução de sistemas lineares}

A cada iteração é necessário resolver
\[
H_i p_i = - \nabla f_i.
\]
Mesmo usando métodos numéricos eficientes, o custo típico é $O(n^3)$ em tempo
para inversão direta ou fatorações, o que se torna proibitivo em alta dimensão.

\subsection*{Dificuldade em obter derivadas de segunda ordem}

Para funções complicadas, a expressão analítica de $\nabla^2 f(x)$ pode ser:
\begin{itemize}
    \item difícil de derivar manualmente;
    \item sujeita a erros de implementação;
    \item custosa de avaliar numericamente.
\end{itemize}
Técnicas como diferenciação automática ajudam, mas ainda há custo de memória e
tempo consideráveis.

\subsection*{Sensibilidade e não convergência em funções não quadráticas}

Em funções não quadráticas:
\begin{itemize}
    \item a Hessiana pode não ser definida positiva em toda a região de interesse;
    \item o método pode divergir se o ponto inicial não estiver suficientemente
          próximo do mínimo;
    \item o método pode convergir para pontos de sela ou máximos se não houver
          controle adicional (por exemplo, via modificações da Hessiana ou
          line search adequado).
\end{itemize}

\section{Por que o Método de Newton Pode Divergir?}

Apesar de suas excelentes propriedades teóricas --- incluindo convergência quadrática 
e obtenção do minimizador em uma única iteração no caso quadrático --- o método de 
Newton pode divergir de forma severa quando aplicado a funções não quadráticas, 
quando o ponto inicial está longe do ótimo ou quando a Hessiana não é definida 
positiva. Nesta seção discutimos, de maneira detalhada, as causas mais comuns de 
divergência.

\subsection*{1. A Hessiana pode não ser definida positiva}

O passo de Newton é obtido resolvendo
\[
H_i p_i = -\nabla f_i,
\]
onde $H_i = \nabla^2 f(x_i)$.

Se a Hessiana for indefinida (isto é, contiver autovalores positivos e negativos), 
então a direção $p_i$ pode apontar para uma direção de \emph{subida}, não de 
descida. Mais ainda:

\begin{itemize}
    \item se $H_i$ tem autovalores negativos, Newton pode mover o ponto em direção 
    a máximos locais;
    \item se $H_i$ é indefinida, Newton mistura direções de subida e descida,
    gerando comportamento altamente instável.
\end{itemize}

\noindent
\textbf{Exemplo:} A função sela
\[
f(x,y) = x^2 - y^2
\]
possui Hessiana
\[
H = \begin{bmatrix}
2 & 0\\
0 & -2
\end{bmatrix},
\]
com um autovalor positivo e um negativo. O método de Newton leva a quedas ao longo 
da direção $x$ e disparos divergentes ao longo da direção $y$. Assim, $x_i$ 
tende ao mínimo enquanto $y_i$ diverge para $\pm\infty$, tornando o método 
praticamente inútil nesse caso.

\subsection*{2. A aproximação quadrática pode ser ruim longe do ótimo}

Newton supõe que $f$ é bem aproximada localmente por uma função quadrática:
\[
f(x) \approx f(x_i) + \nabla f_i^\top (x - x_i) + \tfrac12 (x-x_i)^\top H_i (x-x_i).
\]

Quando $x_i$ está longe do ótimo, esta aproximação pode ser extremamente ruim,
fazendo com que o passo de Newton seja exagerado. O resultado pode ser:

\begin{itemize}
    \item saltos grandes que saem da região de interesse;
    \item oscilações entre regiões distantes;
    \item \emph{overshooting}, em que o método ``passa direto'' pelo vale;
    \item divergência numérica.
\end{itemize}

\noindent
\textbf{Exemplo:} A função $f(x)=\sqrt{x}$ (definida em $x>0$) produz passos de Newton 
que estouram rapidamente e podem sair do domínio da função.

\subsection*{3. O método pode convergir para pontos de sela ou máximos}

Newton procura resolver $\nabla f = 0$, mas não distingue entre tipos de pontos
críticos. Assim, ele pode facilmente convergir para:

\begin{itemize}
    \item máximos locais (onde $H$ é definida negativa),
    \item pontos de sela (onde $H$ é indefinida),
    \item pontos estacionários degenerados.
\end{itemize}

\noindent
\textbf{Exemplo ilustrativo:}  
\[
f(x) = x^3.
\]
Seu único ponto crítico é $x=0$, que é um ponto de sela. No entanto, Newton gera:
\[
x_{i+1} = x_i - \frac{3x_i^2}{6x_i} = \frac{x_i}{2},
\]
convergindo para $0$, interpretando-o erroneamente como solução aceitável.

\subsection*{4. Segunda derivada pequena pode gerar passos catastróficos}

Em uma dimensão, Newton usa:
\[
x_{i+1} = x_i - \frac{f'(x_i)}{f''(x_i)}.
\]

Se $f''(x_i)$ é muito pequeno, então o denominador torna o passo extremamente 
grande, lançando o método muito longe do ótimo. Esse é um dos mecanismos mais 
comuns de divergência em funções suaves, mas relativamente planas.

\subsection*{5. O próximo ponto pode sair do domínio da função}

Quando a função é restrita a um domínio especial (por exemplo, $x>0$ em 
$f(x)=\ln x$), Newton pode gerar $x_{i+1}$ fora do domínio e tornar o processo 
inviável.

\noindent
\textbf{Exemplo:}
\[
f(x) = \ln(x), \qquad x > 0.
\]
Newton produz:
\[
x_{i+1} = 2x_i,
\]
que pode facilmente sair da região em que a função é bem comportada, levando a 
valores inválidos ou divergentes.

\subsection*{6. Sensibilidade extrema ao ponto inicial}

Newton só é garantidamente eficiente quando $x_1$ está suficientemente próximo do 
ótimo. Se o ponto inicial estiver longe:

\begin{itemize}
    \item a aproximação quadrática é incorreta;
    \item a Hessiana pode ter sinal errado;
    \item as derivações sucessivas podem magnificar erros.
\end{itemize}

Esse comportamento explica por que, na prática, usa-se:

\begin{itemize}
    \item line search (Newton amortecido),
    \item métodos de região de confiança,
    \item Hessianas modificadas para garantir positividade.
\end{itemize}

\subsection*{Resumo}

O método de Newton pode divergir devido a:
\begin{itemize}
    \item Hessianas indefinidas ou negativas,
    \item más aproximações quadráticas longe do ótimo,
    \item passos muito grandes quando $f''(x)$ é pequeno,
    \item convergência para pontos de sela,
    \item saída do domínio da função,
    \item alta sensibilidade ao ponto inicial.
\end{itemize}

Quando utilizado sem salvaguardas, o método pode apresentar comportamento
oscilatório ou mesmo explodir numericamente. A inclusão de uma busca linear ou de 
uma região de confiança corrige a maior parte desses problemas, tornando o método 
de Newton mais estável e aplicável em larga escala.



\subsection*{Resumo}

O método de Newton é muito eficiente quando:
\begin{itemize}
    \item $f$ é bem aproximada localmente por uma função quadrática;
    \item a Hessiana é definida positiva e não muda radicalmente;
    \item o custo de computar e manipular a Hessiana é aceitável.
\end{itemize}

Em problemas de grande escala ou com funções muito não lineares, essas condições
raramente são satisfeitas, motivando variantes como métodos quasi-Newton
(BFGS, L-BFGS) e métodos baseados apenas em primeira ordem.

\section{Conclusion}

Newton's method uses second-order information to build a local quadratic model
of the objective function and update the iterate by solving a linear system
involving the Hessian. For quadratic functions, it reaches the minimum in a single
iteration. For general nonlinear functions, it can exhibit fast (quadratic)
convergence near a local minimizer, but its practical use is limited by:
\begin{itemize}
    \item the cost of computing and storing the Hessian;
    \item the need to solve large linear systems;
    \item potential divergence or convergence to saddle points without safeguards.
\end{itemize}

These limitations motivate many modern variants that trade off exact second-order
information for cheaper approximations, preserving some of Newton's advantages
while remaining computationally feasible.

\end{document}
